\documentclass[11pt]{article}
\usepackage[T1]{fontenc}
\usepackage[utf8]{inputenc}
\usepackage{graphicx}
\usepackage{amsmath,amssymb,amsthm}
\usepackage{booktabs}
\usepackage{tabularx}
\usepackage{geometry}
\geometry{a4paper,margin=2.6cm}
\usepackage[hidelinks]{hyperref}
\usepackage{enumitem}
\usepackage{xcolor}

\newtheorem{definition}{Definition}
\newtheorem{proposition}{Proposition}
\newtheorem{remark}{Remark}

% Flags for items still to be settled before circulation.
\newcommand{\tbd}[1]{\textcolor{red}{\textbf{[TBD: #1]}}}
% Marks figures derived in this appendix rather than taken from a source.
\newcommand{\derived}{\textsuperscript{$\dagger$}}

\title{Aviation taxation proposals\\[0.4em]
\large Appendix to the \emph{Report on alternative options for EU own resources}}
\author{}
\date{July 2026}

\begin{document}

\maketitle

\begin{abstract}
\noindent This appendix sets out three nested proposals for aviation taxation. The
first is a package of tax \emph{instruments} --- a tax on vacant seats, a carbon
tax, and a progressive premium-cabin levy --- whose rates can be justified on
efficiency, environmental and distributional grounds respectively. The second is
a \emph{conditional unilateral policy} that applies any subset of those
instruments while earmarking a majority of revenue to global funds, and that
automatically halves its rates on flights connecting to another country
operating the same policy. The third is a \emph{coalition mechanism} that
formalises the norm the conditional policy embodies: that extraterritorial
taxation of flights is legitimate when, and only when, a majority of the revenue
is directed to global public goods through a representative allocation
procedure.

\medskip
\noindent Symbols: figures marked \derived{} are derived in this appendix from
published data and are not taken directly from a cited source. Items marked
\tbd{\ldots} require input before circulation.
\end{abstract}

\tableofcontents

\section{Introduction}
\label{sec:intro}

This appendix argues for three things, each building on the last:

\begin{enumerate}[label=\arabic*)]
  \item a set of aviation tax reform proposals (Section~\ref{sec:instruments});
  \item a unilateral tax policy, using any component from 1), that earmarks at
  least a certain fraction (e.g.\ 60\%) of revenues to global funds for global
  public goods or global redistribution, and that automatically divides the tax
  rates by two for any flight connecting to another country already having such
  a policy in place (Section~\ref{sec:conditional});
  \item a solidarity levy coalition mechanism that builds on 2) and formalises
  the norm it supports: coalition members would respect the norm that unilateral
  extraterritorial taxation of flights is legitimate only if at least a certain
  fraction (e.g.\ 60\%) of revenue is earmarked to global funds whose
  eligibility for receiving funding is determined by a representative vote
  (e.g.\ population-weighted qualified majority voting at the UN). Moreover,
  coalition members would implement the conditional policy defined by 2),
  thereby incentivising other countries to join
  (Section~\ref{sec:coalition}).
\end{enumerate}

The EU could implement 2) and initiate 3). Simply implementing unilateral
extraterritorial carbon pricing for flights with full revenue retention, as it
did in 2012, would risk triggering opposition by other countries that will
likely again consider it a breach of existing international norms. By instead
implementing 2) and initiating 3), the EU could establish a legitimate
architecture that would not only address the externalities associated with
flights but also create a stable source of funding for truly additional global
public good provision and global redistribution.

\subsection{Why this is timely}
\label{subsec:timely}

Three developments make the question live rather than hypothetical.

\emph{First, the EU has just re-opened the extraterritoriality question.} On
17~July 2026 the Commission tabled its revision of the EU Emissions Trading
System (COM(2026)~616). Its Article~28b(2) assessment found that states
participating in CORSIA account for less than 70\% of international aviation
emissions --- below the threshold that would have legally obliged the Commission
to extend the ETS to all departing international flights. Rather than trigger
the full extension, the proposal keeps ``stop-the-clock'' and instead extends
the scope from 2029 to flights departing the EEA and landing in third countries
within 5{,}000~km of the geographic centre of the EU, for four years and subject
to a 2032 review, together with all incoming and departing business-jet flights.
ICAO issued a statement of concern the same day, describing the proposal as
duplicative of CORSIA. The EU has therefore chosen a deliberately restrained
scope precisely to manage the legitimacy problem this appendix addresses --- and
has paid for that restraint in coverage (Section~\ref{subsec:scope}).

\emph{Second, the multilateral norm has hardened against earmarking in
particular.} The 42nd ICAO Assembly (2025) reaffirmed CORSIA as the only global
market-based measure applying to CO\textsubscript{2} from international
aviation, and called on states to raise concerns about duplicative measures
``including aviation-related taxes and levies, particularly where such measures
are intended to mobilize climate finance for sectors beyond international
aviation.'' This is important, and uncomfortable for the argument made here: the
ICAO position objects \emph{more} strongly, not less, where revenue leaves the
aviation sector. Section~\ref{sec:legal} confronts this directly. The short
answer is that the relevant legitimacy forum is contested, and that part of the
purpose of proposal 3) is to shift it.

\emph{Third, a coalition of willing states already exists.} At the Fourth
International Conference on Financing for Development (Seville, July 2025),
eight states --- France, Kenya, Barbados, Spain, Somalia, Benin, Sierra Leone
and Antigua \& Barbuda --- launched the Premium Flyers Solidarity Coalition
under the Seville Platform for Action, with EU support and with the Global
Solidarity Levies Task Force secretariat. That coalition is a natural nucleus
for proposal 3), and its existing commitments (introduce a premium levy;
harmonise upwards; tax private jets) are a subset of what
Section~\ref{sec:coalition} proposes.

\emph{Fourth, and specific to this report.} The Own Resources Decision tabled in
July 2025, COM(2025)~574, proposes five new own resources with an estimated
yield of about €58~billion per year, and the European Parliament has set a floor
of €60~billion per year while explicitly inviting consideration of alternatives
beyond the Commission's basket. Aviation is a candidate the Commission's basket
does not include, other than indirectly through ETS1 revenues. Whether, and in
what form, an earmarked aviation levy can serve as an own resource at all is a
question this appendix must answer rather than assume; see
Section~\ref{subsec:ownresources}.

\subsection{The proposed instruments at a glance}

\begin{table}[htbp]
\centering
\small
\caption{Proposed instruments, the role each plays, and indicative rates.
Rates are indicative and, where marked, derived in this appendix; they are not
a fitted optimum for the EU.}
\label{tab:instruments}
\begin{tabularx}{\textwidth}{@{}p{3.1cm} X p{4.5cm}@{}}
\toprule
\textbf{Instrument} & \textbf{Role} & \textbf{Indicative rate} \\
\midrule
Tax on vacant seats
& Achieve the socially efficient load factor. Corrects the
\emph{utilisation margin}: a dynamically pricing airline leaves more seats empty
than is socially optimal.
& 50\% of the average ticket price on the route, per vacant seat
\newline \emph{(Dörpinghaus et al.)} \\
\addlinespace
Carbon tax
& (i) Internalise climate damages;
\newline (ii) correct the sector's low effective consumption-tax wedge relative
to the rest of the economy;
\newline (iii) act on the \emph{flight-supply margin}, which the vacancy tax
alone cannot target.
& (i) \$190--230/tCO\textsubscript{2} (EPA 2023 central SC-CO\textsubscript{2},
2020 emissions to 2030, 2020~USD);
\newline (ii) see Section~\ref{subsec:wedge} --- \emph{no defensible single
number yet};
\newline (iii) $\approx$\$100--115/tCO\textsubscript{2}\derived{} at the
passenger-neutral calibration \\
\addlinespace
Progressive premium-cabin and private-jet levy
& Compensate for the limited progressivity of the overall tax system, and reach
globally mobile high earners whom no single income tax reaches.
& \$165 per premium ticket (with \$82.50 economy) or \$360 per premium ticket
alone, each raising \$100bn globally \emph{(ICCT 2025)}; private-jet fuel excise
per GSLTF model rates \\
\bottomrule
\end{tabularx}
\end{table}

Two features of Table~\ref{tab:instruments} deserve emphasis at the outset.

First, the instruments target \emph{different margins} and are therefore
complements, not substitutes. The vacancy tax acts on how fully a scheduled
aircraft is filled; the carbon tax acts on how many aircraft are scheduled at
all. Using one instrument to do both jobs --- as a conventional ad valorem
ticket tax does --- necessarily sacrifices one margin to the other. This is the
central analytical claim of Dörpinghaus et al.\ and the reason the package is
worth more than the sum of a debate about ticket-tax rates.

Second, the carbon tax rate is doing three conceptually distinct jobs in
Table~\ref{tab:instruments}, and the three do not coincide. Only role~(i) has a
well-established number. Role~(iii) is a \emph{constrained} rate: it is whatever
carbon price holds passenger numbers at their status quo level given the vacancy
tax, and it happens to come out at roughly half the EPA central social cost of
carbon. Role~(ii) is theoretically the most delicate and, as
Section~\ref{subsec:wedge} explains, we do not think a number should be put in
that cell without further work. Presenting a single ``optimal carbon tax''
figure would obscure this.

\section{Tax reform proposals}
\label{sec:instruments}

We consider three sets of tax reforms that we argue would be desirable in terms
of their effects on global welfare. Ideally all three would be implemented
together, each at its respective optimal ambition level. However, each is
valuable on its own and could be implemented as such. We begin with the reform
that comes closest to avoiding creating losers, and which should therefore be
easiest to build support for.

\subsection{Tax reforms to transport the same number of people with fewer and
fuller planes}
\label{subsec:pfp}

We propose the following tax reform as a minimal reform that could garner broad
international support.

\begin{definition}[Passenger-flow-preserving efficiency-maximising reform]
\label{def:pfp}
The reform consists of introducing the following tax instruments:
\begin{enumerate}[label=\arabic*)]
  \item a tax on vacant seats, whereby airlines pay for each vacant seat at
  departure a certain percentage of the average ticket price on the route;
  \item a carbon tax.
\end{enumerate}
The tax on vacant seats is set so that the socially efficient vacancy rate is
reached, and the carbon tax is adjusted such that the number of passengers
transported remains unchanged as compared with the status quo. Where an ad
valorem ticket tax is already in place, it is abolished as part of the package.
\end{definition}

\begin{remark}
The reform is defined by two \emph{targets} (efficient utilisation; unchanged
passenger volume), not by two rates. This matters politically: the package can
be described without committing in advance to a carbon price, and the
passenger-volume target is verifiable ex post. It also matters analytically: the
vacancy tax rate that achieves the first target turns out to be robust to the
level at which the second is set (see Section~\ref{subsec:results},
point~\ref{item:robustrate}).
\end{remark}

\subsubsection{Results from the underlying model}
\label{subsec:results}

Dörpinghaus et al.\ provide a full microeconomic model of this reform. The model
places a route in an infinite horizon with an endogenous number of flights:
customers arrive according to a Poisson process, draw a valuation and a
preferred departure date, and choose among currently selling flights by
comparing valuation net of schedule-disutility against posted prices. Each
flight is operated by a distinct airline choosing a Markovian pricing policy,
attention is restricted to symmetric Markov Perfect Equilibria that are
asymptotically stable under damped best-response dynamics, and flight frequency
is pinned down by a zero-profit free-entry condition. The model is calibrated to
a status quo vacancy rate of 20\%, a marginal cost of filling a seat equal to
20\% of the average ticket price, and a current ad valorem tax burden of
6.5\%.

The model yields the following results for the passenger-flow-preserving reform.

\begin{enumerate}
  \item The tax on vacant seats is approximately \textbf{50\% of the average
  ticket price}.
  \item\label{item:robustrate} This rate is approximately unchanged if a
  different target for the number of flights is chosen instead of the status quo
  --- the vacancy tax is, in this sense, a target-independent instrument for the
  utilisation margin.
  \item The vacancy rate falls from \textbf{20\% to about 6\%}. (In the simpler
  fixed-frequency variants of the model the optimum is lower still: about 3\%
  under constant-elasticity demand and about 1.5\% under exponential demand.
  The 6\% figure is the one to quote, since it comes from the full model with
  endogenous entry.)
  \item Flight frequency falls by \textbf{17.8\%} (from $M=3.00$ to $M=2.47$ in
  the model's units).
  \item Emissions fall by approximately \textbf{17.8\%} where applied, since in
  this model emissions are determined almost entirely by the number of flights
  and are affected only slightly by the vacancy rate. Applied globally this is
  of the order of \textbf{150\,MtCO\textsubscript{2}}\derived{} per year against
  a global aviation total of roughly 850--950\,Mt; applied to all flights with
  at least one EU endpoint, roughly \textbf{0.29} of that
  (Section~\ref{subsec:scope}).
  \item Aviation tax revenue rises by \textbf{\$34.8~billion} globally in
  partial equilibrium, and by \textbf{\$55.1~billion}\derived{} once the
  consumption taxes and corporate income taxes accruing in the rest of the
  economy from redirected spending are included
  (Section~\ref{subsec:ge}). At EU scope: about \textbf{\$10~billion} and
  \textbf{\$16~billion}\derived{} respectively.
  \item Global economic surplus rises by \textbf{\$72~billion} once pre-existing
  wedges in the rest of the economy are accounted for (\$34.8~billion in partial
  equilibrium). At EU scope, about \textbf{\$21~billion}\derived{}.
\end{enumerate}

For orientation, Table~\ref{tab:endpoints} reproduces the model's main
counterfactuals. It is worth showing in full, because the single number one
quotes depends entirely on which constraint one imposes, and a reader who
encounters two different welfare figures without this table will reasonably
suspect an error.

\begin{table}[htbp]
\centering
\small
\caption{Policy endpoints in the free-entry model of Dörpinghaus et al.
Welfare gains in the final column are partial-equilibrium and exclude the
general-equilibrium correction of Section~\ref{subsec:ge}. $M$ is the number of
flights simultaneously selling.}
\label{tab:endpoints}
\begin{tabular}{@{}lrrrrr@{}}
\toprule
Endpoint & $M$ & Pax & Revenue & Carbon rev. & Fiscal dividend \\
 & & [bn] & [bn USD] & [bn USD] & [bn USD] \\
\midrule
Status quo (6.5\% ticket tax)   & 3.00 & 5.20 & 1000.0 & 0.0 & 0.0 \\
Vacancy-tax optimum, $M$ fixed  & 3.00 & 6.30 &  931.6 & 0.0 & $-41.1$ \\
Zero-profit $M^\star$           & 2.77 & 5.83 &  923.7 & 0.0 & $-44.4$ \\
\textbf{Passenger-neutral}      & 2.47 & 5.20 &  904.5 & 82.1 & $+34.8$ \\
Social-surplus-neutral          & 2.29 & 4.83 &  885.4 & 121.7 & $+73.1$ \\
\midrule
\multicolumn{6}{@{}l@{}}{\emph{Partial-equilibrium welfare gain [bn USD]:}
0.0 \quad 106.6 \quad 80.5 \quad \textbf{34.8} \quad 0.0} \\
\bottomrule
\end{tabular}
\end{table}

Three readings of Table~\ref{tab:endpoints} matter for policy.

\begin{itemize}
  \item The \emph{largest} welfare gain (\$106.6~billion) comes from the vacancy
  tax alone with flight frequency held fixed. But holding frequency fixed
  requires compensating airlines for lost operating profit, and it costs the
  exchequer \$41.1~billion relative to the ticket tax it replaces. It is the
  right number for the question ``how large is the utilisation distortion?'' and
  the wrong number for a policy proposal.
  \item The \emph{passenger-neutral} endpoint is the one to lead with. It raises
  revenue, raises surplus, cuts emissions, and leaves passenger numbers
  unchanged --- four properties that are individually attractive to different
  constituencies and jointly rare.
  \item The \emph{social-surplus-neutral} endpoint shows the fiscal headroom:
  \$73.1~billion globally, at the cost of 7.1\% fewer passengers and no net
  surplus gain. It is the natural upper bound of the ``no welfare loss'' range
  and a useful reference point for a revenue-driven discussion.
\end{itemize}

\subsubsection{The implied carbon price}
\label{subsec:implied}

The passenger-neutral endpoint is defined by a target, so the carbon price that
implements it is an output rather than an input. It can be recovered
approximately from the model's carbon revenue.

At that endpoint the model raises \$82.1~billion in carbon revenue while
reducing flights, and hence approximately emissions, by 17.8\%. Applying the
reform globally to an aviation sector emitting $E$ tonnes at the status quo
gives post-reform emissions of $0.822\,E$ and therefore an implied carbon price
of $\$82.1\text{bn}/(0.822\,E)$. For $E$ between 0.88 and 1.00~GtCO\textsubscript{2}
--- ATAG reports 882\,Mt for 2023 and IATA a higher figure for 2024 --- this
gives:

\begin{center}
\begin{tabular}{@{}lr@{}}
\toprule
Status quo global aviation CO\textsubscript{2} & Implied carbon price\derived{} \\
\midrule
0.88\,Gt & \$114/tCO\textsubscript{2} \\
0.95\,Gt & \$105/tCO\textsubscript{2} \\
1.00\,Gt & \$100/tCO\textsubscript{2} \\
\bottomrule
\end{tabular}
\end{center}

So the passenger-flow-preserving package implies a carbon price of roughly
\textbf{\$100--115/tCO\textsubscript{2}} --- on the order of \emph{half} the EPA
central social cost of carbon of \$190/t for 2020 emissions rising to \$230/t by
2030 (2020~USD, 2\% near-term Ramsey discount rate).

This is an important and easily missed point. The passenger-flow-preserving
reform is \emph{not} the fully internalising reform. It is a deliberately
constrained package that buys efficiency gains and revenue without reducing
passenger volumes, and it therefore leaves roughly half the Pigouvian gap
unclosed. That constraint is what makes it politically attractive; it should not
be presented as sufficient on climate grounds. Section~\ref{subsec:more-carbon}
takes up what going further would involve.

\subsubsection{From global to EU scope: the coverage ladder}
\label{subsec:scope}

The global figures above must be scaled to the geographic scope actually
applied. This is not a detail: the multiplier varies by a factor of nearly three
across the options currently on the table, and using a single number invites
challenge.

The share of global aviation CO\textsubscript{2} covered can be derived as
follows. T\&E report that flights departing European airports emitted 195\,Mt in
2025 and accounted for more than 23\% of worldwide aviation emissions, implying
a global total of about 848\,Mt; EU-departing flights accounted for 154\,Mt of
this, i.e.\ 18.2\%. EASA reports that 61\% of EU27+EFTA departure emissions are
on flights with a destination outside EU27+EFTA. Assuming approximate symmetry
between departures and arrivals gives:

\begin{center}
\begin{tabular}{@{}p{8.4cm}rr@{}}
\toprule
Scope & MtCO\textsubscript{2} & Share of global\derived{} \\
\midrule
Intra-EU only & 60 & 0.07 \\
Commission proposal from 2029 (intra-EEA $+$ departures $\leq$5{,}000\,km) & $\approx$90 & $\approx$0.11 \\
All flights departing the EU & 154 & 0.18 \\
\textbf{All flights with at least one EU endpoint} & \textbf{248} & \textbf{0.29} \\
\bottomrule
\end{tabular}
\end{center}

The 0.29 figure used throughout this appendix therefore corresponds to the
\emph{most} extensive scope --- taxing flights in both directions --- which is
precisely the scope the Commission declined to adopt on 17~July 2026. The scope
actually proposed covers roughly 0.11 of global aviation emissions, i.e.\ about
\emph{a third} of what the 0.29 multiplier assumes.

\begin{remark}[Two reasons the scaling is not exactly linear]
First, revenue and welfare do not scale with emissions coverage in strict
proportion, because the reform changes behaviour and because EU routes differ
from the global average in stage length, load factor and premium share.
Second, and more importantly, a unilaterally applied reform loses effectiveness
to \emph{hub leakage}: passengers reroute via untaxed third-country hubs. This
is not speculative --- the Commission's stated rationale for the 5{,}000\,km
radius is precisely that it captures Istanbul, Dubai and Doha and so addresses
hub leakage. A vacancy tax may be more exposed to this than a carbon tax, since
it can be avoided by restructuring an itinerary into segments that fall outside
the taxing jurisdiction. The 0.29 (or 0.11) multiplier should therefore be read
as an \emph{upper bound} on the emissions effect at EU scope.
\end{remark}

\subsubsection{The general-equilibrium correction}
\label{subsec:ge}

The model's direct estimate of the welfare gain from the passenger-neutral
reform is \$34.8~billion. This is a partial-equilibrium figure: it implicitly
treats the rest of the economy as a numeraire good entering utility
quasilinearly, i.e.\ as having a zero wedge between consumer price and marginal
cost. In fact the rest of the economy has substantial wedges, and the reform
releases real resources into it.

The correction runs as follows. Consumption taxes average 11.5\% globally and
market-power-induced mark-ups a further 14.5\%, giving an economy-wide wedge of
$1.115 \times 1.145 - 1 = 0.28$. At the status quo, 80\% of airline costs are
fixed, and by the zero-profit condition total costs equal revenue net of taxes,
i.e.\ $\$(1000-64.8)$~billion. A 17.8\% reduction in flight frequency therefore
saves
\[
0.80 \times 0.178 \times \$935.2\text{bn} = \$133.2\text{bn}
\]
in real resources, which are redirected to other sectors where each dollar of
cost yields \$1.28 of consumer benefit. The partial-equilibrium calculation
therefore omitted $0.28 \times \$133.2\text{bn} = \$37.4\text{bn}$, giving a
total welfare gain of
\[
\$34.8\text{bn} + \$37.4\text{bn} = \$72.2\text{bn}.
\]

The same redirection generates additional revenue outside aviation:
$0.115 \times \$133.2\text{bn} = \$15.3\text{bn}$ in consumption taxes and
$0.26 \times 0.145 \times \$133.2\text{bn} = \$5.0\text{bn}$ in corporate income
tax on the mark-ups, for a total fiscal effect of
$\$34.8 + \$15.3 + \$5.0 = \$55.1\text{bn}$\derived{}.

\begin{remark}[Three cautions on this correction]
\emph{(i)} It is arithmetically sensitive: using the unrounded wedge of 0.2767
rather than 0.28 gives \$36.8bn and a total of \$71.6bn rather than \$72.2bn.
The figure should be quoted as ``about \$72~billion''.
\emph{(ii)} It applies only where flights are reduced. There is no analogous
correction to the \$106.6~billion vacancy-tax-optimum row of
Table~\ref{tab:endpoints}, because frequency is held fixed there and no
resources are released. This asymmetry is correct but confusing, and should be
stated whenever both numbers appear.
\emph{(iii)} It assumes the released resources are in fact redeployed to sectors
bearing the average wedge. Where flight reductions occur at slot-constrained
airports, the freed capacity is a scarce right that will be reallocated to other
flights rather than released to the wider economy, and the correction does not
apply. Roughly \tbd{share of EU departures at Level~3 coordinated airports}
of EU departures are at fully coordinated airports, so this qualification is not
marginal.
\end{remark}

\subsubsection{Incidence}
\label{subsec:incidence}

In general, aviation tax reforms raise the difficult question of how the
incidence is split between travellers and the hospitality sector.
Conceptually it is clear that this is determined by the relative elasticities of
the demand for flights and the supply of hospitality services (see
\href{https://meetings.imf.org/en/-/media/files/publications/staff-climate-notes/2024/english/clnea2024003.pdf}{IMF
Staff Climate Note CLNEA/2024/003}). However, estimating these elasticities is
difficult, and as a result there appear to be no studies giving confident
estimates.

An attractive feature of the passenger-flow-preserving reforms is that, at least
in the simplest economic models, they create no burden for the hospitality
sector: the number of passengers transported is unchanged by construction.
Moreover, the impact on travellers --- via changes in prices, changes in
availability of flights and changes in flight frequency --- is well captured by
the model of Dörpinghaus et al.\ and is thus already taken into account in the
estimates above.

Two qualifications should be stated explicitly rather than left to critics.

First, passenger neutrality holds \emph{in aggregate}, not
\emph{destination by destination}. The carbon component of the package falls in
proportion to emissions and therefore bears disproportionately on long-haul
routes. Since the aggregate passenger constraint is met by short-haul volumes
expanding as long-haul contracts, tourism-dependent long-haul destinations lose
even when global passenger numbers are held fixed. This is not a peripheral
objection: small island developing states are among the most exposed, and
several --- Barbados, Antigua \& Barbuda, Fiji, the Marshall Islands --- are
participants in the very solidarity levy coalition on which
Section~\ref{sec:coalition} builds. Any serious version of this proposal needs a
route-level or region-level incidence analysis, and probably a compensating
allocation rule. \tbd{route-level incidence analysis, or an explicit commitment
that the global fund allocation formula weights tourism-dependence}

Second, the ``no burden on hospitality'' claim relies on the reform being
applied at a scope wide enough that rerouting is not an option. At EU scope,
part of the burden is shifted to competing hubs rather than absorbed, which
helps EU-bound tourism but at the cost of the emissions objective.

\subsubsection{Robustness and what the model does not capture}
\label{subsec:robustness}

The estimates above rest on a model, and the model's authors are careful about
its limits. A policy appendix that presents the 50\% figure as settled will not
survive technical scrutiny. The following caveats should travel with the
numbers.

\begin{itemize}
  \item \textbf{The result is demand-specification dependent.} For sufficiently
  thick-tailed distributions of valuations, a tax on vacant seats \emph{reduces}
  welfare: pushing prices down early causes seats to be sold to low-valuation
  travellers before high-valuation travellers arrive. The favourable result
  holds for the constant-elasticity and exponential families that dominate the
  literature, and the authors conjecture --- but do not prove --- that welfare
  is single-peaked in the vacancy tax.

  \item \textbf{Time-varying demand weakens the instrument.} A constant vacancy
  tax acts on the load-factor margin and the intertemporal-allocation margin
  simultaneously and cannot separate them. Under time-varying demand it is a
  second-best instrument only, and the calibrated optimal vacancy rates and
  welfare gains are lower than in the stationary case.

  \item \textbf{The calibration target is now stale.} The model is calibrated to
  a 20\% vacancy rate (80\% load factor). ATAG reports average occupancy of
  83.6\% for 2025, and European load factors are higher still. Recalibrating to
  current utilisation would mechanically shrink the estimated gains, since the
  distortion being corrected is smaller than assumed.
  \tbd{recalibrate to a 16\% vacancy rate and report the revised headline
  figures}

  \item \textbf{Vacancy is not pure waste.} The model treats an empty seat at
  departure as a lost opportunity. In practice airlines hold inventory
  deliberately: for connecting passengers under origin-and-destination controls,
  as a buffer against no-shows and operational disruption, and because
  overbooking carries denied-boarding liabilities. Under Regulation
  261/2004 those liabilities are €250--600 per passenger. Driving load factors
  from 80\% to 94\% would increase denied boardings substantially, and the
  associated compensation and disruption costs are absent from the welfare
  calculation. This is both a real welfare cost and, in the EU specifically, a
  legal-interaction problem.
  \tbd{estimate the EU261 interaction; it may be material at a 6\% vacancy rate}

  \item \textbf{Network and market structure are stylised.} Each flight is
  operated by a distinct airline --- a deliberate competitive limit, chosen to
  show that the distortion does not vanish with competition, but one that
  abstracts from network revenue management, hub-and-spoke structure, belly
  cargo, and slot scarcity. Where slots bind, the free-entry zero-profit
  condition that pins down flight frequency fails, and with it the mechanism
  generating the \$133.2~billion resource saving.

  \item \textbf{The underlying paper is not final.} Several quantities in the
  manuscript, including the carbon tax rate at the passenger-neutral endpoint,
  are still blank. The figures in this appendix should be checked against the
  final version before circulation.
  \tbd{confirm all figures against the final Dörpinghaus et al.\ manuscript}
\end{itemize}

None of this undermines the qualitative case, which is robust across every
specification the authors examine: airlines leave too many seats empty, and a
tax on vacant seats is the natural instrument. It does mean the quantitative
claims should be presented as an order of magnitude with a stated calibration,
not as a point estimate.

\subsection{Additional carbon taxation}
\label{subsec:more-carbon}

The passenger-flow-preserving reform of Section~\ref{subsec:pfp} involves an
increase in carbon taxation to roughly \$100--115/tCO\textsubscript{2}, but by
construction only as much as leaves passenger numbers unchanged. There are two
distinct arguments for going further.

\subsubsection{Internalising climate damages}

The first is straightforward Pigouvian logic. The EPA's 2023 central estimates
put the social cost of CO\textsubscript{2} at \$190 per tonne for 2020
emissions, rising to \$230 by 2030 and \$308 by 2050 (2020~USD, 2\% near-term
Ramsey discount rate); the range across the 2.5\%, 2.0\% and 1.5\% discount
rates spans roughly \$140 to \$380 for 2030 emissions. Two cautions apply when
citing these figures in an EU document. First, Executive Order 14154 of
January 2025 disbanded the US Interagency Working Group, withdrew its estimates
and directed EPA to consider removing the social cost of carbon from federal
decision-making, so ``EPA'' is no longer an uncontested authority for the
number; the underlying peer-reviewed literature is the more durable citation.
Second, the figures are in 2020 dollars and for a specified emissions year, so a
single ``\$210'' quotation needs its vintage and deflator stated.

Aviation's climate impact also exceeds its CO\textsubscript{2} impact.
Accounting for non-CO\textsubscript{2} effects --- principally
contrail cirrus --- T\&E puts the EU aviation climate impact at 262\,MtCO\textsubscript{2}e
against 154\,Mt of CO\textsubscript{2}, a multiplier of about 1.7. The
uncertainty here is genuine and EASA is explicit about it, but a
CO\textsubscript{2}-only carbon tax is, on current best estimates, pricing
substantially less than half the damage. \tbd{decide whether to propose a
CO\textsubscript{2}e basis or a CO\textsubscript{2} basis with a stated
multiplier}

\subsubsection{The wedge argument, and why we do not put a number on it}
\label{subsec:wedge}

The second argument is that aviation is undertaxed relative to other consumption
even setting externalities aside. International aviation is exempt from VAT and,
under the Energy Taxation Directive and Article~24 of the Chicago Convention,
kerosene is effectively untaxed: in the G20 the average price of kerosene in
2021 was €9 per tonne of CO\textsubscript{2}, against €79 for diesel and €68
for petrol. The natural inference is that the sector's tax wedge should be
raised to the economy-wide average of 28\% computed in
Section~\ref{subsec:ge}, implying an additional ad valorem charge of roughly
20~percentage points on top of the current 6.5\%.

We think that inference is too quick, and that the corresponding cell of
Table~\ref{tab:instruments} should be left open rather than filled with a
plausible-looking number.

The reason is that the standard wedge-equalisation prescription compares
consumer price to \emph{marginal} cost. In aviation, 80\% of the cost of a
flight is fixed at the point of departure and the marginal cost of filling a
seat is around 20\% of the average ticket price. With free entry and zero
profits, price equals average cost, which necessarily and substantially exceeds
marginal cost --- the ratio implied by the calibration is on the order of five.
Applied naively, ``equalise the wedge over marginal cost'' would therefore
prescribe a \emph{subsidy}, not a tax, which is clearly not the intended
conclusion. The mark-up over marginal cost in a free-entry industry with large
fixed costs is not a measure of market power to be taxed away; it is the
mechanism by which fixed costs are recovered.

What the general-equilibrium correction in Section~\ref{subsec:ge} does is
different and defensible: it values resources \emph{released} by the sector at
the wedge prevailing where they are redeployed. That is a statement about the
welfare accounting of a given reform, not a prescription for a target wedge in
aviation. Deriving the correct second-best ad valorem charge for a
free-entry industry with fixed costs, VAT exemption and an untaxed input is a
genuine open question, and one on which this appendix should either do the work
or decline to give a number. We recommend the latter for now.
\tbd{either derive the second-best ad valorem rate for a fixed-cost free-entry
sector, or drop role (ii) from Table~\ref{tab:instruments} and rest the case on
roles (i) and (iii)}

\begin{remark}
There is a cleaner version of the undertaxation argument that avoids the
marginal-cost problem entirely: the VAT exemption on international tickets is an
anomaly relative to the treatment of domestically consumed services, and
removing an exemption is not the same as computing an optimal wedge. That
argument supports adding a charge approximately equal to the VAT that would
otherwise apply, and it does not require any claim about mark-ups. It is
narrower but far more robust, and it is probably the one to make in a report on
own resources.
\end{remark}

\subsection{Progressive taxes on tickets}
\label{subsec:progressive}

The third component addresses the limited progressivity of the overall tax
system. See
\href{https://theicct.org/wp-content/uploads/2025/03/ID-302-\%E2\%80\%93-Aviation-levy_report_final-v2.pdf}{ICCT
(2025), ``Designing an equitable aviation climate levy'' (ID-302)}, which
compares six designs: an air passenger duty, an aviation fuel tax, a frequent
flying levy, an air miles levy, a luxury aviation levy on premium cabins, and a
ticket levy with rebates.

\subsubsection{What the designs deliver}

The distributional differences between designs are large. A frequent flying levy
and an air miles levy concentrate 80\% and 84\% of total levy charges
respectively on high-income households; an air miles levy shifts about
\$30~billion, or 30\% of the total burden, from low- and middle-income to
high-income households relative to a flat duty. ICCT concludes that a frequent
flying levy or air miles levy is most suitable for raising funds for global
climate finance, and that a ticket levy with rebates achieves similar
distributional effects while being easier to implement. Instruments that exempt
infrequent flying minimise the impact on once-a-year non-business trips.

On rates: to raise \$100~billion globally from premium cabins alone would
require about \$360 per first- or business-class ticket; a design with a
9$\times$ multiplier on a flat ticket tax --- \$165 per premium ticket and
\$82.50 per economy ticket --- raises the same amount with a lower headline
premium rate. CE~Delft estimates that a global levy on premium flyers could
raise about €78~billion per year. Across the fuel, ticket and frequent flyer
designs, roughly two-thirds of revenue is raised in high-income countries and
about a quarter in upper-middle-income countries, which is the property that
makes these instruments attractive as a source for global funds rather than
merely as domestic revenue.

\subsubsection{The theoretical basis, stated honestly}

A commodity tax on premium cabins needs a justification, because the
Atkinson--Stiglitz theorem implies that with weakly separable preferences and an
optimal non-linear income tax, differential commodity taxation is
welfare-reducing. Three arguments are available, and they are not equally
strong.

\begin{enumerate}[label=(\alph*)]
  \item \emph{The income tax is not optimal.} This is true but proves too much;
  it licenses almost any differential commodity tax and gives no guidance on
  rates.
  \item \emph{Preferences are not weakly separable in the relevant sense.} If
  premium travel is a status or positional good, consumption generates a
  negative externality on other consumers and a corrective tax is justified
  independently of the income tax. This is a real argument but requires an
  estimate of the status component to pin down a rate.
  \item \emph{The relevant tax base is not reachable by any income tax.} This is
  the strongest argument in the international context and the one specific to
  the present proposal. A premium-cabin levy at EU scope falls substantially on
  non-residents --- globally mobile high earners whose income no single
  jurisdiction taxes effectively. A levy on the consumption of internationally
  mobile capital income is not a substitute for a progressive income tax that
  exists; it is a substitute for one that does not.
\end{enumerate}

Argument (c) is the one to lead with in a report on own resources, because it
identifies a base the EU can reach and Member States individually cannot ---
which is precisely the criterion for an own resource rather than a national tax.

\subsubsection{Interaction with the other two instruments}

A premium-cabin levy sits awkwardly alongside the vacancy tax and should not be
introduced without thinking about the interaction. Premium cabins have
structurally lower load factors than economy, so a vacancy tax applied per seat
falls disproportionately on premium capacity --- which is either a feature (it
compounds the progressivity objective) or a problem (it may induce airlines to
reconfigure cabins toward economy, raising seat counts, lowering yields, and
partially undoing the emissions reduction per passenger). The direction of the
net effect is not obvious and depends on the relative elasticities.
\tbd{model the cabin-configuration response to a per-seat vacancy tax}

\section{Conditional aviation tax policies with earmarking for global funds}
\label{sec:conditional}

We propose to implement the following at the EU level.

\subsection{The earmarking obligation}

\begin{definition}[Earmarking obligation]
\label{def:earmark}
Let $r \in [0,1]$ be a \emph{retention rate}. A jurisdiction respects the
\emph{earmarking obligation at retention rate $r$} with respect to a stream of
aviation tax revenue if it
\begin{enumerate}[label=(\roman*)]
  \item retains at most the proportion $r$ of that revenue for its own budget;
  and
  \item transfers the remaining proportion $1-r$ to global funds currently
  recognised as eligible under the procedure of
  Definition~\ref{def:eligibility}, allocated across eligible funds in
  proportion to the aggregate receipts of each fund from all contributing
  jurisdictions.
\end{enumerate}
\end{definition}

The proportionality requirement in (ii) is not incidental. Without it, a
jurisdiction could satisfy the letter of the obligation while directing its
entire contribution to whichever eligible fund most closely served its own
foreign-policy interests, which would hollow out the claim that the revenue
finances \emph{globally} chosen priorities. Tying each contributor's allocation
to the aggregate pattern of receipts means that the composition of global
spending is determined collectively even though contributions are made
unilaterally. The reference value for $r$ in this appendix is $r = 0.4$, i.e.\
60\% earmarked, but Section~\ref{subsec:participation} shows that $r$ is the
single most consequential parameter in the design and that 0.4 may be too
demanding.

\subsection{The conditional levy}

\begin{definition}[Conditional aviation solidarity levy at rate $\tau$]
\label{def:conditional}
Consider the following sequence of policy components, each to be applied subject
to the earmarking obligation of Definition~\ref{def:earmark}:

\begin{description}[leftmargin=2.2cm,style=nextline]
  \item[$P_0$:] Tax all domestic flights at rate $\tau$.
  \item[$P_1$:] Tax all international flights at rate $\tfrac{\tau}{2}$.
  \item[$P_2$:] For all international flights coming from or going to a country
  not doing $P_0$ and $P_1$, increase the rate to $\tau$.
  \item[$P_3$:] For all international flights coming from or going to a country
  not doing $P_2$, increase the rate to $\tau$.
  \item[$\;\;\vdots$] \
  \item[$P_n$:] For all international flights coming from or going to a country
  not doing $P_{n-1}$, increase the rate to $\tau$.
\end{description}

Define $P^{*} := \bigcup_{k=0}^{\infty} P_k$.
\end{definition}

\begin{remark}[Two corrections to the original formulation]
The union must begin at $k=0$ rather than $k=1$; otherwise domestic flights are
never taxed and $P_2$ has no antecedent to refer to. Also, the recursion is
well-founded --- $P_k$ is defined purely in terms of observation of other
countries' compliance with $P_{k-1}$, which is itself already defined --- so
there is no circularity, and because every escalation step raises the rate to
$\tau$ and never beyond, the union is a well-defined policy with rates bounded by
$\tau$.
\end{remark}

\begin{remark}[Rate-matching for partial compliance]
For ease of exposition Definition~\ref{def:conditional} treats compliance as
binary. In practice a more general version is preferable: if a country operates
$P_0$ and $P_1$ at some lower rate $\tau' < \tau$, then $P_2$ would increase the
rate only to $\tau - \tfrac{\tau'}{2}$ rather than to the full $\tau$. This is
analogous to the way the EU's CBAM credits carbon prices already paid abroad, and
it removes the cliff-edge that a binary rule creates.
\end{remark}

\subsubsection{What \texorpdfstring{$P^{*}$}{P*} actually amounts to}

The infinite hierarchy in Definition~\ref{def:conditional} is what makes the
policy \emph{norm-complete}: it penalises not only countries that fail to levy,
but countries that fail to penalise countries that fail to levy, and so on. This
closes off second-order free riding. But the hierarchy also makes the policy look
forbiddingly complex, and an administration asked to implement ``an infinite
union of policies'' will reasonably object. The following observation, which we
believe is worth stating explicitly in any presentation of the mechanism, shows
that the complexity is entirely in the justification and not at all in the
implementation.

\begin{proposition}[Fixed-point characterisation]
\label{prop:fixedpoint}
A country operates $P^{*}$ if and only if it
\begin{enumerate}[label=(\alph*)]
  \item taxes all domestic flights at rate $\tau$;
  \item taxes an international flight to or from country $j$ at rate
  $\tfrac{\tau}{2}$ if $j$ operates $P^{*}$, and at rate $\tau$ otherwise; and
  \item respects the earmarking obligation on all of the above.
\end{enumerate}
\end{proposition}

\begin{proof}[Sketch]
Suppose $i$ operates $P^{*}$ and consider a flight between $i$ and $j$.
If $j$ operates $P^{*}$ then $j$ operates $P_0$ and $P_1$, so the escalation in
$P_2$ is not triggered; and $j$ operates $P_{n-1}$ for every $n \geq 3$, so no
subsequent escalation is triggered either. The rate remains $\tfrac{\tau}{2}$.
If $j$ does not operate $P^{*}$, let $k$ be least with $j$ failing $P_k$. If
$k \in \{0,1\}$ the escalation in $P_2$ applies; if $k \geq 2$ the escalation in
$P_{k+1}$ applies. Either way the rate is $\tau$. Conversely, a country
satisfying (a)--(c) satisfies every $P_k$ by construction.
\end{proof}

So implementing $P^{*}$ requires a tax authority to observe one binary fact
about each partner country --- whether it has enacted the same rule --- and
nothing more. No regress needs to be computed. The hierarchy is the
\emph{argument} for why the rule is closed under its own logic; the rule itself
is a single conditional.

\begin{remark}[On uniqueness]
Proposition~\ref{prop:fixedpoint} identifies $P^{*}$ as a fixed point of the
operator that maps a candidate rule to the rule ``halve for those who follow the
candidate, full rate otherwise''. The union construction reaches this fixed point
by iteration from the non-escalating base, which makes it the natural or least
such fixed point. We have not established that it is the unique one, and the
claim should not be made without proof. Note, however, that the multiplicity
that matters practically is not in the definition of the rule but in the
\emph{set of countries that adopt it}: the adoption game has the tipping
structure described in Section~\ref{subsec:participation} and can be expected to
have multiple equilibria. That is a feature of any club mechanism and is
precisely why an explicit coordinating institution (Section~\ref{sec:coalition})
adds value over unilateral adoption.
\end{remark}

\subsubsection{Route-tax invariance and why joining is attractive}
\label{subsec:participation}

The halving rule has a property that is not obvious from the definition and that
does most of the work in generating participation incentives.

\begin{proposition}[Route-tax invariance]
\label{prop:invariance}
Let $i$ operate $P^{*}$. Then the total tax borne by a flight between $i$ and
$j$ equals $\tau$ whether or not $j$ operates $P^{*}$. What changes is only the
\emph{division} of that revenue: $j$ collects nothing if it stays out, and half
if it joins.
\end{proposition}

\begin{proof}
If $j$ stays out, $i$ levies $\tau$ and $j$ levies nothing. If $j$ joins, each
levies $\tfrac{\tau}{2}$.
\end{proof}

This is a materially different incentive structure from that of a conventional
border-adjustment club. In a CBAM-style club, joining means imposing a new cost
on one's own economy in exchange for avoiding a border charge. Here, on routes to
existing members, joining imposes \emph{no additional burden at all} on the
joiner's travellers --- the burden is already being levied by the member at the
other end --- and simply transfers half of it into the joiner's own fiscal
control. Relative to the counterfactual of staying out, that component of joining
is a pure revenue gain.

The cost of joining lies elsewhere: a member must also tax its own domestic
flights at $\tau$, and its flights to and from \emph{non}-members at $\tau$,
where previously it taxed neither. A stylised threshold follows. Write $B_{jj}$
for $j$'s domestic base, $B_{jk}$ for the base on the route between $j$ and $k$,
and $M$ for the current membership. Valuing retained revenue at par and
travellers' burden at par, and ignoring deadweight loss, $j$ gains from joining
whenever
\[
\underbrace{\frac{r\tau}{2}\sum_{i \in M} B_{ij}}_{\text{revenue captured at no new burden}}
\;>\;
\underbrace{(1-r)\,\tau\Big(B_{jj} + \sum_{k \notin M} B_{jk}\Big)}_{\text{revenue exported on newly taxed routes}},
\]
that is, whenever
\[
\beta_j \;:=\; \frac{\tfrac{1}{2}\sum_{i \in M} B_{ij}}{B_{jj} + \sum_{k \notin M} B_{jk}}
\;>\; \frac{1-r}{r}.
\]

Two conclusions follow, and both bear directly on the choice of the 60\% figure.

\emph{First, there is a tipping structure.} As $M$ grows, the numerator of
$\beta_j$ rises and the denominator falls, so the incentive to join strengthens
with every accession. Membership is a strategic complement, which is what makes
an initial coalition of sufficient size the critical variable.

\emph{Second, the retention rate is the binding design parameter.} At $r = 0.4$
(60\% earmarked) the condition is $\beta_j > 1.5$: a country joins only once the
halved base on member routes exceeds one and a half times its combined domestic
and non-member base. That is demanding, and for most countries it will not be
satisfied until the coalition is already very large. At $r = 0.6$ (40\%
earmarked) the condition relaxes to $\beta_j > 0.667$, which is satisfiable by a
much wider set of countries at much smaller coalition sizes.

There is therefore a genuine trade-off between the ambition of the earmarking
norm and the size of the coalition that will sustain it, and the 60\% figure sits
on the demanding side of it. We would recommend either lowering the earmarked
share, or --- better --- making it \emph{increasing in coalition size}, so that
the norm tightens automatically as the club grows and the participation
constraint slackens. A schedule such as ``40\% earmarked until members account
for a quarter of global departures, rising to 60\% thereafter'' would preserve
the ambition while removing the early-stage barrier.
\tbd{decide whether to propose a fixed or coalition-size-contingent earmarking
share; this is the most consequential open design choice in the appendix}

\begin{remark}[What the stylised threshold omits]
The condition above values a euro of retained revenue equally with a euro of
traveller burden, which amounts to assuming a marginal cost of public funds of
one and no distributional weighting. It also ignores deadweight loss, the
domestic political salience of taxing one's own citizens' domestic flights
relative to taxing foreigners' international flights (which is likely to be
substantial and to work against joining), and any value the joining country
places on the global public goods financed. The last of these works in favour of
joining and is the reason a mechanism of this kind can outperform what the
narrow fiscal calculus suggests. The threshold should be read as organising the
trade-off, not as a calibrated prediction.
\end{remark}

\subsection{The revenue base for the earmarking obligation}
\label{subsec:base}

The earmarking obligation must be applied to a specified base, and the choice of
base changes the fiscal arithmetic more than any other implementation detail. It
is worth setting out explicitly, because the natural reading of the proposal
turns out to leave Member State treasuries worse off.

Take the passenger-flow-preserving package of Section~\ref{subsec:pfp} applied at
a scope covering 0.29 of global aviation. Scaling the global figures:

\begin{table}[htbp]
\centering
\small
\caption{Fiscal arithmetic at EU scope (coverage share 0.29), \$bn per year.
All figures derived\derived{} by scaling the global results of
Table~\ref{tab:endpoints}.}
\label{tab:eu-fiscal}
\begin{tabular}{@{}lr@{}}
\toprule
Gross new revenue (carbon \$82.1bn + vacancy \$17.5bn, scaled) & 28.9 \\
Pre-existing ad valorem ticket tax, abolished by the package & $-18.8$ \\
\midrule
\textbf{Net new public revenue (fiscal dividend)} & \textbf{10.1} \\
\bottomrule
\end{tabular}
\end{table}

Now compare two readings of the earmarking obligation.

\begin{description}[leftmargin=1.2cm]
  \item[(A) 60\% of \emph{gross} new revenue.] \$17.3bn is transferred to global
  funds and \$11.6bn retained. But the package abolished \$18.8bn of
  pre-existing ticket-tax revenue. Public budgets in the EU therefore end up
  \$7.2bn \emph{worse off} than the status quo, even though the reform generates
  \$10.1bn of net new revenue. The \$7.2bn is not financed out of the dividend;
  it is a redirection of pre-existing national revenue to global funds.
  \item[(B) 60\% of the \emph{net fiscal dividend}.] \$6.1bn is transferred to
  global funds and \$4.0bn retained on top of a treasury position held whole. No
  pre-existing revenue is redirected.
\end{description}

Reading (A) is very likely fatal in the Council, which must adopt an Own
Resources Decision by unanimity. Reading (B) is fiscally coherent but yields a
global-fund contribution roughly a third the size, and the appendix should own
that difference rather than quote the larger number while implying the
non-redirection of existing revenue.

A third possibility is conceptually the most attractive. What requires
international legitimation is not the levy as a whole but its
\emph{extraterritorial} component --- the taxation of flights to and from third
countries. Revenue raised on intra-EU flights raises no legitimacy question at
all. This suggests applying the earmarking obligation only to revenue raised on
the extraterritorial base, which at EU scope is roughly 76\% of the
total.\footnote{Derived from the coverage decomposition in
Section~\ref{subsec:scope}: intra-EU 60\,Mt of a 248\,Mt EU-endpoint total.} Doing
so ties the size of the obligation directly to the scope of the extraterritorial
claim being made, which is exactly the logic the coalition norm is meant to
express.

\begin{description}[leftmargin=1.2cm]
  \item[(C) 60\% of the extraterritorially-raised portion of the net dividend.]
  Roughly \$4.6bn to global funds, with the remainder retained and treasuries
  held whole.
\end{description}

We recommend (C) as the principled option and (B) as the simpler fallback, and we
recommend against (A). \tbd{choose between (B) and (C) and re-derive all
headline global-fund figures accordingly}

\begin{remark}[Own resources, not national revenue]
Table~\ref{tab:eu-fiscal} speaks of ``public budgets'' deliberately. The
pre-existing ticket taxes are \emph{national}; a new EU-level levy would be an
\emph{EU} resource, and Member States would be compensated through reduced
GNI-based contributions in the usual way. The reform therefore involves a
vertical shift of revenue as well as a net increase, and the two should be
presented separately: Member States will read any table that nets them together
as a loss of fiscal sovereignty.
\end{remark}

\section{International aviation solidarity levy coalition mechanism}
\label{sec:coalition}

Building on the conditional policies of Section~\ref{sec:conditional}, the EU
could initiate the following international mechanism. The definition below is our
proposed content for what was previously a placeholder; each clause is separable
and the numbered parameters are all negotiable.

\begin{definition}[Aviation solidarity levy coalition mechanism]
\label{def:mechanism}
The mechanism consists of the following rules.

\medskip
\noindent\textbf{R1 (Substantive commitment).} Each member implements the
conditional aviation solidarity levy $P^{*}$ of
Definition~\ref{def:conditional} at a rate no lower than a floor $\underline\tau$,
using any combination of the instruments of Section~\ref{sec:instruments},
subject to the earmarking obligation at a retention rate no higher than $r$.

\medskip
\noindent\textbf{R2 (Instrument neutrality).} Members are free to choose among a
carbon tax, a tax on vacant seats, a distance-based passenger duty, a
premium-cabin levy, or any combination, provided the aggregate burden meets the
floor $\underline\tau$ on a common measurement basis. Compliance is assessed on
revenue raised per unit of the agreed base, not on instrument choice.

\medskip
\noindent\textbf{R3 (The legitimacy norm).} Members affirm that unilateral
extraterritorial taxation of international flights is legitimate if and only if
the earmarking obligation is respected in respect of the extraterritorially
raised revenue. Members undertake (i) not to challenge, retaliate against, or
support challenges to, measures of other members that comply with the norm; and
(ii) to treat non-compliant extraterritorial taxation --- including their own
past measures --- as falling outside the norm.

\medskip
\noindent\textbf{R4 (Eligibility of recipient funds).} See
Definition~\ref{def:eligibility}.

\medskip
\noindent\textbf{R5 (Allocation).} Earmarked revenue is allocated across
eligible funds in proportion to their aggregate receipts from all members, as in
Definition~\ref{def:earmark}(ii). No member may condition its contribution to an
eligible fund on that fund's decisions.

\medskip
\noindent\textbf{R6 (Verification).} Members publish annually the base, rates,
revenue raised, and amounts transferred, disaggregated by domestic,
intra-coalition and extra-coalition flights. Compliance is assessed by an
independent secretariat against a published methodology. Verification determines
only whether a member is treated as operating $P^{*}$ for the purposes of other
members' halving rules; it carries no other sanction.

\medskip
\noindent\textbf{R7 (Entry into force).} The mechanism enters into force when
members jointly account for at least a share $\underline{s}$ of global
departures, measured on the coverage basis of Section~\ref{subsec:scope}. Before
that threshold, members may implement $P^{*}$ unilaterally but the
non-retaliation commitments of R3 do not bind.

\medskip
\noindent\textbf{R8 (Accession and exit).} Accession is open to any state
enacting R1 and accepting R3 and R6. Exit requires notice of at least
\tbd{notice period}; on exit, other members' halving rules cease to apply to the
exiting state at the end of the notice period.

\medskip
\noindent\textbf{R9 (Review).} The parameters $\underline\tau$, $r$ and
$\underline{s}$ are reviewed at fixed intervals, with $r$ following the
coalition-size-contingent schedule if one is adopted
(Section~\ref{subsec:participation}).
\end{definition}

\begin{definition}[Eligibility of recipient funds]
\label{def:eligibility}
A global fund is eligible to receive earmarked revenue if it is designated as
such by a population-weighted qualified majority vote of UN member states, held
at fixed intervals, under a published set of criteria comprising at least:
additionality of the financing relative to existing commitments; provision of a
global public good or global redistribution; independent audit; and open
publication of allocations.
\end{definition}

\subsection{Why the eligibility vote is the load-bearing element}

R4 and Definition~\ref{def:eligibility} are doing more work than their length
suggests, and it is worth being explicit about why.

The central objection to any extraterritorial levy is that the taxing state is
appropriating a base that is not exclusively its own. Earmarking answers this
only if the recipients are not chosen by the taxing state. If the EU both taxes
extraterritorially \emph{and} decides where the money goes, the earmarking is
cosmetic and the objection stands with full force. What makes the architecture
defensible is that eligibility is determined by a procedure in which the EU is a
small minority: a population-weighted vote of UN members gives the EU roughly
5--6\% of the weight. The EU is thereby committing to raise revenue whose
destination it does not control.

This is a real commitment and should be presented as such rather than
minimised. It is also the element most likely to be resisted internally, for
exactly the same reason it is persuasive externally.

Two design questions follow, and neither has an obvious answer.

First, \emph{population-weighted voting is not a settled institution and does not
currently exist at the UN.} The General Assembly votes one state one vote;
population weighting would have to be created for this purpose. Specifying it as
a novel procedure is defensible in a research appendix but will read as
speculative in a policy report unless an existing or near-existing forum is
identified. Candidates worth considering are a weighted vote within an existing
fund's governing body, the UN Framework Convention's arrangements, or a plurilateral
body established by the coalition itself with weighted membership open to all
states. \tbd{identify a concrete forum, or state explicitly that the mechanism
requires creating one}

Second, \emph{population weighting is not obviously the right weighting.} It
gives India and China jointly about 35\% of the weight and gives small island
developing states --- the constituency with the strongest claim on aviation
levy revenue and, as noted in Section~\ref{subsec:incidence}, the one most
exposed to the levy's incidence --- almost none. A defence of population
weighting on representativeness grounds has to confront the fact that the
resulting allocation may be less favourable to the most vulnerable than a
one-state-one-vote allocation would be. A hybrid --- population-weighted with a
floor per state, or a double majority of states and population --- may be
easier to defend on both grounds.

\subsection{Relationship to the Premium Flyers Solidarity Coalition}

The mechanism above should be positioned as a deepening of an existing
initiative rather than a competing one. The Premium Flyers Solidarity Coalition
launched at Seville in July 2025 by France, Kenya, Barbados, Spain, Somalia,
Benin, Sierra Leone and Antigua \& Barbuda already commits its members to
introduce a premium-cabin levy where none exists, to harmonise existing levies
upwards, and to tax private jets, with proceeds directed to resilience and fair
transitions. Three of the four elements of R1 are therefore already
politically live among a group that includes two EU Member States.

What Definition~\ref{def:mechanism} adds to that coalition is precisely what it
currently lacks: the \emph{conditionality} that generates participation
incentives (R1 via $P^{*}$), the \emph{norm} that makes extraterritorial
application defensible (R3), and a \emph{collective procedure} for determining
recipients (R4). The Task Force's own legal handbook already supplies drafting
guidance and model rates for the premium-cabin and private-jet instruments,
which reduces the implementation burden considerably. The natural sequencing is
therefore for the EU to adopt the conditional structure and offer it to the
existing coalition, rather than to propose a new institution.

\section{Legal and institutional feasibility}
\label{sec:legal}

\subsection{International aviation law}

The legal position differs sharply across the three instruments, and the
appendix should not treat them as a single package for this purpose.

\emph{Distance-based passenger duties} are the safe case. Larsson et al.\ (2019)
conclude that only distance-based air passenger taxes are unambiguously legal
under international law. Nothing in the Chicago Convention prohibits them:
Article~24 exempts fuel already on board from customs duties, and Article~15
prohibits charges imposed solely in respect of the right of transit, entry or
exit, neither of which reaches a tax on the sale of a ticket. It is precisely for
this reason that the United Kingdom retained its Air Passenger Duty and abandoned
its 2010 attempt to replace it with a per-plane tax, notwithstanding that the
government appeared persuaded that a per-plane tax could raise welfare by
increasing load factors (Dresner 2010).

\emph{Carbon pricing of international flights} is contested but not clearly
unlawful. The 2012 episode is instructive: the EU exempted extra-EEA flights
from November 2012 under threat of retaliatory measures, not following any
adverse legal determination. The 42nd ICAO Assembly's reaffirmation of CORSIA
exclusivity is a resolution, not a treaty obligation, and Assembly resolutions do
not create binding law. The obstacle is political and reputational rather than
strictly legal --- which is not to say it is weaker.

\emph{A tax on vacant seats} is the novel case and carries the most legal
uncertainty. It is a charge on the operator rather than on the passenger, and it
has no precedent. There is a plausible argument that it is structurally a ticket
tax on a notional passenger and inherits the ticket tax's legality; there is a
contrary argument that a per-aircraft charge assessed at departure is closer to
the per-plane tax the UK abandoned, and closer to the Article~15 prohibition on
charges in respect of exit. This matters because the vacancy tax is the
distinctive contribution of the underlying research and the appendix's most
original recommendation.

Two mitigations are worth building into the design. First, express the vacancy
tax as a per-seat charge indexed to distance, so that it resembles as closely as
possible the instrument that is known to be lawful. Second, note that its
revenue yield is modest (\$17.5bn globally, against \$82.1bn from the carbon
component), so that if the legal risk proves unacceptable the package can be
implemented without it at a limited fiscal cost --- though at the cost of the
efficiency gain that motivates the whole exercise.
\tbd{commission a legal opinion on the vacancy tax under the Chicago Convention
and a sample of bilateral Air Services Agreements}

Bilateral Air Services Agreements are a further constraint that the appendix
currently does not address at all. Many contain reciprocal tax exemption
provisions, and the EU's competence is limited where Member State bilaterals
apply. A survey of the relevant clauses is a prerequisite for any concrete
proposal. \tbd{ASA survey}

\subsection{The ICAO norm problem}

Section~\ref{subsec:timely} noted that the 42nd Assembly objected specifically to
levies ``intended to mobilize climate finance for sectors beyond international
aviation''. This deserves to be met directly rather than passed over, because it
appears to invert the appendix's central premise: the document argues that
earmarking to global funds \emph{buys} legitimacy, while the revealed ICAO
position is that earmarking outside the sector \emph{costs} legitimacy.

Three responses are available, and we think they should all be made.

First, the two claims are about different audiences. ICAO's position reflects a
forum in which industry interests are strongly represented and in which states
with large hub carriers are influential. It is not a proxy for the view of the
UN membership at large, for whom revenue directed to global funds through a
representative procedure is the attractive feature rather than the objectionable
one. Part of the purpose of R3 and R4 is to shift the forum in which the
legitimacy of aviation taxation is adjudicated.

Second, the ICAO objection is strongest against measures that are
\emph{both} extraterritorial \emph{and} unilaterally allocated. The mechanism
proposed here concedes allocation to a collective procedure, which is a
substantive answer to the double-charging and appropriation concerns rather than
a rhetorical one.

Third, the objection has less force against instruments that are not
market-based measures at all. A premium-cabin levy is not a carbon price and does
not duplicate CORSIA on any reading; the ICAO Assembly's concern about
double-charging simply does not apply to it. This is a further reason to lead the
external case with the progressive instrument of
Section~\ref{subsec:progressive} rather than with carbon pricing, even though
carbon pricing carries the larger revenue and the stronger efficiency rationale.

\subsection{EU law and the own resources question}
\label{subsec:ownresources}

Three EU-specific constraints bear on the proposal, and one of them may be
decisive for its form.

\emph{Unanimity.} A tax measure under Article~113 TFEU and an Own Resources
Decision under Article~311 TFEU both require unanimity in Council, with the
Parliament only consulted in the latter case. This is the binding political
constraint and the reason Section~\ref{sec:implementation} treats the ETS route
as a serious alternative rather than a fallback.

\emph{Budgetary universality.} EU budget revenue is in principle not assigned to
specific expenditure. A design in which 60\% of a levy never enters the EU budget
at all, or enters it earmarked to particular external recipients, sits uneasily
with that principle and with the Parliament's budgetary prerogatives. There is a
straightforward architectural fix: route the earmarked share as EU
\emph{expenditure} --- a global public goods facility within the MFF --- rather
than as revenue diverted before it reaches the budget. This preserves
universality and parliamentary control, and it has the political advantage that
the EU can point to a visible budget line. It has the disadvantage that
expenditure is renegotiated every MFF cycle, which weakens exactly the
predictability that makes the revenue attractive to recipient funds. The
trade-off between commitment and constitutional propriety is real.
\tbd{legal service view on whether an assigned-revenue or expenditure route is
available}

\emph{Own resource criteria.} For the purposes of this report the relevant test
is not only revenue yield but EU value added --- whether the base is one that
Member States cannot effectively tax individually. Here the proposal scores
genuinely well, and the argument should be made explicitly. A premium-cabin levy
falls substantially on non-residents whose income no Member State taxes; a
levy applied at EU scope is far less vulnerable to hub leakage than the same levy
applied by one Member State, since intra-EU rerouting is neutralised; and the
conditional structure of $P^{*}$ only functions at a scale that makes the halving
rule attractive to third countries. Each of these is an argument that this
revenue \emph{belongs} at EU level rather than merely that it would be
convenient there.

\section{Implementation options}
\label{sec:implementation}

We have proposed three tax instruments that could be implemented at EU level if
unanimity can be reached among Member States. For the carbon tax, there is an
alternative: implementation through ETS1, which has the significant advantage of
not requiring unanimity.

If unanimity could be reached, it might nonetheless be advantageous to implement
an additional EU-wide carbon tax on flights alongside ETS1 rather than instead of
it. Given the undertaxation of flights even setting aside externalities
(Section~\ref{subsec:wedge}), there is an advantage in having a separate
instrument whose rate can be set on aviation-specific grounds, rather than having
the carbon price on flights determined residually by the ETS1 price, which is set
by conditions across the whole covered economy.

The 17~July 2026 ETS proposal changes the practical calculus in three ways that
the appendix should reflect.

\begin{enumerate}
  \item \emph{The scope question is now open in legislation.} The Commission has
  proposed the 5{,}000\,km radius from 2029 with a review in 2032, and has said
  explicitly that it may consider extending to all departing flights if CORSIA
  does not deliver. The conditional structure of $P^{*}$ is directly relevant to
  that 2032 decision: it offers a way to extend scope while reducing rather than
  increasing the risk of retaliation. This is the natural point of legislative
  entry.
  \item \emph{The proposal already contains a deduction mechanism} to avoid
  double carbon pricing where CORSIA and the ETS both apply. That mechanism is
  structurally the same idea as the rate-matching remark following
  Definition~\ref{def:conditional} --- crediting action taken elsewhere --- and
  the precedent is useful.
  \item \emph{Business jets are brought into scope without a distance limit.}
  This is the one place where the Commission has accepted unlimited
  extraterritorial reach, presumably because the political economy of private
  aviation is unusually favourable. It suggests that the private-jet component of
  the progressive instrument is the easiest place to start.
\end{enumerate}

A staged sequence follows naturally, and we would recommend presenting it in
this order rather than as a single package:

\begin{enumerate}[label=\textbf{Stage \arabic*.},leftmargin=1.8cm]
  \item Private-jet fuel excise and premium-cabin levy at EU scope, with the
  earmarking obligation applied to the extraterritorial portion. Legally
  safest, politically most popular (Section~\ref{sec:support}), no CORSIA
  duplication, and already the subject of an existing coalition.
  \item The conditional structure $P^{*}$ applied to that levy, establishing the
  norm and the halving incentive at low stakes.
  \item Carbon pricing via ETS1 scope extension at the 2032 review, with the
  conditional structure and earmarking already established, so that scope
  extension arrives as an application of an accepted norm rather than as a fresh
  unilateral act.
  \item The vacancy tax, subject to the legal opinion and to recalibration
  against current load factors.
\end{enumerate}

\section{Public support}
\label{sec:support}

The distributional design of the instruments in
Section~\ref{subsec:progressive} is supported by public opinion to an unusual
degree. Research by Oxfam and Greenpeace finds that three out of four people
across thirteen countries believe wealthier air travellers should pay more tax in
view of their outsized climate impact. The underlying inequality is stark:
private aviation emissions rose 46\% between 2019 and 2023, and high-income
countries generated about 70\% of aviation CO\textsubscript{2} between 1980 and
2019.

Two features of the political economy should temper any inference from these
figures.

First, the opposition is concentrated and well organised while the support is
diffuse. IATA dismissed the Premium Flyers Solidarity Coalition declaration in
November 2025, arguing that the proposal lacked credible analysis of its effects
on the countries it aims to support and that it duplicates CORSIA. The
industry's argument that CORSIA already channels climate finance to developing
economies through credit purchases --- with expected investment of more than
\$100bn between 2024 and 2035 --- is one the appendix should engage with on the
merits, since the additionality of those flows is precisely the question that
Definition~\ref{def:eligibility}'s additionality criterion is designed to
address.

Second, support for taxing \emph{other people's} premium travel is not the same
as support for the domestic component of $P^{*}$, which taxes citizens' own
domestic flights at the full rate $\tau$. The polling evidence supports the
progressive instrument specifically; it should not be cited in support of the
package as a whole. \tbd{polling on the domestic component, and on earmarking to
global funds rather than to domestic budgets --- the latter is likely to be the
harder sell and we have no evidence on it}

\section{Open questions}
\label{sec:open}

The following are the points at which we judge the argument to be currently
incomplete, in rough order of how much they matter.

\begin{enumerate}
  \item \textbf{The retention rate.} Section~\ref{subsec:participation} suggests
  that 60\% earmarking may be too demanding to generate accession at realistic
  coalition sizes. A coalition-size-contingent schedule should be modelled
  properly, with the tipping threshold characterised.
  \item \textbf{The base for the earmarking obligation.} Options (A), (B) and (C)
  of Section~\ref{subsec:base} differ by a factor of nearly four in the resulting
  global-fund contribution. A choice is needed before any figure is quoted.
  \item \textbf{Recalibration.} The 20\% status quo vacancy rate is stale;
  current global occupancy is 83.6\% and European load factors are higher. All
  headline figures should be recomputed.
  \item \textbf{The eligibility forum.} Population-weighted UN voting does not
  currently exist and its distributional properties may not favour the states
  with the strongest claim. A concrete alternative is needed.
  \item \textbf{Route-level incidence.} Aggregate passenger neutrality conceals
  a shift from long-haul to short-haul that falls on tourism-dependent
  developing economies, including coalition members.
  \item \textbf{The wedge argument.} Either derive the second-best ad valorem
  rate for a free-entry sector with large fixed costs, or rest the undertaxation
  case on the narrower and more robust VAT-exemption argument.
  \item \textbf{Legal opinion on the vacancy tax}, and a survey of bilateral Air
  Services Agreement tax clauses.
  \item \textbf{The EU261 interaction} at load factors of 94\%.
  \item \textbf{Uniqueness of the fixed point} in
  Proposition~\ref{prop:fixedpoint}, and characterisation of the equilibrium set
  of the accession game.
\end{enumerate}

\end{document}
